\documentclass[10pt,letterpaper]{article}
\usepackage[lmargin=2.2in,rmargin=0.7in,tmargin=1.0in,bmargin=1.0in]{geometry}
\usepackage{notes}
\begin{document}

\begin{center}
\Large\textbf{3.3 - The Shift Formula}\\
\normalsize (Section 3.1)
\end{center}

% ---------------------------------------
\section{Need More Transform Formulas}

\subsection{What We Can Solve}

\begin{wraptable}{r}{0.32\textwidth}
\vspace{-\baselineskip}
\centering
\renewcommand{\arraystretch}{1.6}
\begin{tabular}{|c|c|c|}
\hline
\multicolumn{3}{|c|}{%
    \shortstack{\textbf{Current Table}\\\textbf{of Transforms}}}\\
\hline
&$x(t)$ & $X(s)$\\
\hline
1&$1$ & $\tfrac{1}{s}$\\
2&$t$ & $\tfrac{1}{s^2}$\\
3&$t^n$ & $\tfrac{n!}{s^{n+1}}$\\
4&$e^{at}$ & $\tfrac{1}{s-a}$\\
5&$\cos(\beta t)$ & $\tfrac{s}{s^2+\beta^2}$\\
6&$\sin(\beta t)$ & $\tfrac{\beta}{s^2+\beta^2}$\\
\hline
\end{tabular}
% \vspace{-\baselineskip}
\end{wraptable}

Our goal is to use Laplace transforms to solve more
types of second-order DEs. Our current table lets us
recover constants, polynomials, exponentials,
sines, and cosines.

If the solution to the DE is a sum of exponentials
$$
    x=c_1e^{\lambda_1t}+c_2 e^{\lambda_2t} \implies X=\frac{As+B}{(s-\lambda_1)(s-\lambda_2)}
$$
then we can use partial fraction decomposition on $X(s)$. Doing so will yield sums of formulas that appear in the right column of the transform table. 

The same is true for purely imaginary eigenvalues $\lambda=\pm \beta i$. The solution is a sum
of sine and cosine terms
$$
    x = c_1\cos\beta t+c_2\sin\beta t \implies X=\frac{As+B}{s^2+\beta^2}
$$
The transforms of these terms also appear in the table. However, other root types produce functions
that our table does not yet include.

\subsection{What We Cannot Solve Yet}

\begin{wraptable}{r}{0.6\textwidth}
\vspace{-\baselineskip}
\renewcommand{\arraystretch}{1.4}
\setlength{\tabcolsep}{5pt}
\begin{tabular}{|l|c|c|}
\hline
\textbf{Root type}
& \textbf{Example}
& \shortstack{\textbf{Can we solve}\\\textbf{with Laplace?}}\\
\hline
Distinct real & $x''+3x'+2x=0$ & $\checkmark$\\
\hline
Purely imaginary & $x''+4x=0$ & $\checkmark$\\
\hline
Repeated real & $x''+4x'+4x=0$ & \\
\hline
Complex ($\alpha\ne0$) & $x''+2x'+5x=0$ & \\
\hline
\end{tabular}
\end{wraptable}
Now consider an equation with a repeated root:
\[
x''+4x'+4x=0.
\]
which has the general solution
\[
x(t)=c_1e^{-2t}+c_2te^{-2t}.
\]
We know the transforms of $t$ and $e^{-2t}$ separately,
but our table does not yet include their product $te^{-2t}$. 
Similarly, consider an equation with complex roots
having a nonzero real part:
\[
x''+2x'+5x=0.
\]
The roots are $\lambda=-1\pm2i$, so we expect
\[
x(t)=e^{-t}\bigl(c_1\cos2t+c_2\sin2t\bigr).
\]
Our \textit{current table of transforms} does not have a formula for the transform of this solution. 

\subsection{Goal}
Notice that both the repeated roots and complex roots cases have a solution of the form 
\[
e^{at}f(t).
\]
This is the missing piece. The goal of today is to learn the \textit{exponential shift property} and solve the 2nd order DEs with repeated or complex roots.

% ---------------------------------------

\section{The Exponential Shift Property}

\subsection{The Shift Property}
Multiplying a function by an exponential produces a shift in the Laplace domain. 
\begin{theorem}[Shift Property]
If
\[
\mathcal L[f(t)]=F(s),
\]
then
\[
\mathcal L[e^{at}f(t)]=F(s-a).
\]
\end{theorem}
This will be given as Formula 7 in the Laplace table.
Although proofs are not usually required, the calculation below illustrates a technique you will use in the homework. We therefore present it as a worked example.

\begin{example}[1]{
Prove the formula 
\[
\mathcal L[e^{at}f(t)]=F(s-a).
\]
}
    
By definition 
$$
\mathcal L[e^{at}f(t)] = \int_0^\infty e^{at} f(t) e^{-st}\,dt.
$$
We can manipulate the exponent and create a new variable $s^* = s-a$:
\begin{align*}
    \mathcal L[e^{at}f(t)] &= \int_0^\infty e^{at} f(t) e^{-st}\,dt&&\com{by definition}\\
    &= \int_0^\infty f(t) e^{-t(s-a)}\,dt\\
    &= \int_0^\infty f(t) e^{-ts^*}\,dt &&\com{let $s^*=s-a$}\\
    &=F(s^*) &&\com{by definition}
\end{align*}
Therefore, we have
$$
\mathcal L[e^{at}f(t)] = F(s^*) = F(s-a).
$$
\end{example}

The shift property lets us build new transform formulas
from ones already in our table. First, identify the
\textit{base transform} $f(t)\to F(s)$.
Then account for the exponential factor $e^{at}$ by
replacing $s$ with $s-a$. Let's apply this to $te^{-2t}$, the term that appeared
in our repeated-root solution.

\begin{example}[2]{
Compute the transform
\[
\mathcal L[te^{-2t}].
\]
}
First, identify the \textit{base transform}. Since our function is $t$ times $e^{-2t}$ the base function is $f(t)=t$. Find the transform of $f(t)$ by itself 
\[
\mathcal{L}[f(t)]=\mathcal L[t]=\frac{1}{s^2} =: F(s).
\]
Since the function contains $e^{-2t}$, the shift is
\[
a=-2 \implies F(s-a)=F(s+2).
\]
Therefore
\[
\mathcal L[te^{-2t}]
= F(s+2)=
\boxed{
\frac{1}{(s+2)^2}.
}
\]
\end{example}

% --------------



\subsection{Using the Shift Property in Reverse}

To solve differential equations, we use the shift
property backward:
\[
\mathcal{L}^{-1}[F(s-a)]=e^{at}f(t).
\]
First, recognize the base function $F(s)$.
Then identify the shift amount $a$. Then apply the formula by substitution.

\begin{example}[3]{
Compute the inverse Laplace transform of 
\[
X(s)=
\frac{s-1}{(s-1)^2+9}.
\]
}
The base transform function and its inverse are
\[
F(s)=\frac{s}{s^2+9},
\qquad
f(t)=\cos3t.
\]
Next, we observe by comparison that the original function can be written as
\[
F(s-1)=\frac{s-1}{(s-1)^2+9}.
\]
Thus, if $F(s-a)=F(s-1)$ then the shift is $a=1$. Finally, we apply the formula and substitute
\[
\mathcal L^{-1}
\left[
\frac{s-1}{(s-1)^2+9}
\right] = e^{at}f(t)=
\boxed{
e^t\cos3t.
}
\]
\end{example}

In Example 3, the numerator and denominator contain
the same shift $s-1$. When they do not match,
we must rewrite the numerator before using the table.

\begin{example}[4]{
Compute the inverse Laplace transform of 
\[
X(s)=
\frac{s}{(s-1)^2+1}
\]
}
First, compare the function to formulas in the table to identify the base transform. We see that this resembles the cosine transformation, and so 
$$
    F(s)=\frac{s}{s^2+1}, \quad f(t)=\cos t.
$$
The denominator suggests a shift $a=1$, but the
shifted cosine formula requires $s-1$ in the numerator.
Write $s=(s-1)+1$ and split the function over addition:
\begin{align*}
\frac{s}{(s-1)^2+1}
&=\frac{(s-1)+1}{(s-1)^2+1}\\
&=\frac{s-1}{(s-1)^2+1} 
+\frac{1}{(s-1)^2+1}.
\end{align*}
Now apply the inverse transform to each term:
\begin{align*}
\mathcal{L}^{-1}\left[\frac{s}{(s-1)^2+1}\right]
&=
\underbrace{\mathcal{L}^{-1}
\left[\frac{s-1}{(s-1)^2+1}\right]}_{\text{shifted cosine}}
+
\underbrace{\mathcal{L}^{-1}
\left[\frac{1}{(s-1)^2+1}\right]}_{\text{shifted sine}}\\
&=e^t\cos t+e^t\sin t.
\end{align*}
Therefore,
\[
\boxed{x(t)=e^t(\cos t+\sin t).}
\]
\end{example}

\textbf{Key idea:} match the numerator to the shift
in the denominator. Any leftover constant produces
a sine term.
% --------------

\section{Solving IVPs}

\subsection{Procedure}
Let's review the general procedure for solving IVPs:
\begin{enumerate}
    \item Transform the entire DE.
    \item Solve for $X$. 
    \item Apply the inverse transform.
\end{enumerate}
The third step is most difficult out of the three. The third step may require a combination of 
\begin{itemize}
    \item partial fraction decomposition,
    \item shift formula, or
    \item completing the square.
\end{itemize}
Here is an example of an IVP with repeated roots.

\begin{example}[5]{
Use Laplace transforms to solve
\[
x''+4x'+4x=0,\qquad x(0)=1,\quad x'(0)=0.
\]
}
\textbf{Step 1: Transform the DE.}
\begin{align*}
\mathcal{L}[x'']+4\mathcal{L}[x']+4\mathcal{L}[x]&=0,\\
(s^2X-sx(0)-x'(0))+4(sX-x(0))+4X&=0,\\
(s^2X-s)+4sX-4+4X&=0.
\end{align*}

\textbf{Step 2: Solve for $X(s)$.}
\[
(s^2+4s+4)X=s+4
\quad\implies\quad
X(s)=\frac{s+4}{s^2+4s+4}.
\]

\textbf{Step 3: Apply the inverse transform.}
The denominator factors as $(s+2)^2$, so
\[
X(s)=\frac{s+4}{(s+2)^2}.
\]
We must deal with the $s$ in the numerator. One option is to partial fraction decomposition
$$
X(s)=\frac{s+4}{(s+2)^2} = \frac{A}{s+2}+\frac{B}{(s+2)^2},
$$
but there is a simpler approach. You can add and subtract 2 in the numerator 
$$
X(s)=\frac{s+4}{(s+2)^2} = \frac{s+4+2-2}{(s+2)^2}=\frac{(s+2)+2}{(s+2)^2} = \frac{1}{s+2}+\frac{2}{(s+2)^2}
$$
The base function in $\frac{2}{(s+2)^2}$ is $F(t)=\frac{1}{s^2}$ which has the inverse $f(t)=t$. The shift is $a=-2$. Therefore, the inverse of the entire function is 
\[
x(t)=\mathcal{L}^{-1}\left[\frac{1}{s+2}+\frac{1}{(s+2)^2}\right]=
\boxed{
e^{-2t}+te^{-2t}.
}
\]
\end{example}

% ------

We now consider an equation with complex roots. Its solution should contain
an exponential multiplied by sine and cosine. 


\begin{example}[6]{
Use Laplace transforms to solve
\[
x''+2x'+5x=0,\qquad x(0)=1,\quad x'(0)=0.
\]
}
\textbf{Step 1: Transform the DE.}
\begin{align*}
\mathcal{L}[x'']+2\mathcal{L}[x']+5\mathcal{L}[x]&=0,\\
(s^2X-sx(0)-x'(0))+2(sX-x(0))+5X&=0,\\
(s^2X-s)+2(sX-1)+5X&=0.
\end{align*}

\textbf{Step 2: Solve for $X(s)$.}
\[
(s^2+2s+5)X=s+2
\quad\implies\quad
X(s)=\frac{s+2}{s^2+2s+5}.
\]

\textbf{Step 3: Apply the inverse transform.}
Even though $X$ is a rational function, partial fraction decomposition will not be effective here. The denominator $s^2+2s+5$ is \textit{irreducible}. 

We want the denominator in the form $(s-a)^2+\beta^2$.
Since $(s+1)^2=s^2+2s+1$, we add and subtract $1$:
\begin{align*}
s^2+2s+5
&=(s^2+2s+1)+5-1\\
&=(s+1)^2+4.
\end{align*}
This technique is called \textbf{completing the square}.
Now
\[
X(s)=\frac{s+2}{(s+1)^2+4}.
\]
The denominator reveals the shift $a=-1$ and frequency
$\beta=2$. To match the shifted cosine formula,
write the numerator as $s+2=(s+1)+1$:
\begin{align*}
X(s)
&=\frac{s+1}{(s+1)^2+4}
+\frac{1}{(s+1)^2+4}\\
&=\frac{s+1}{(s+1)^2+2^2}
+\frac12\frac{2}{(s+1)^2+2^2}.
\end{align*}

\textbf{Step 4: Apply the inverse transform.}
The first term gives a shifted cosine, and the second
gives a shifted sine:
\[
x(t)=e^{-t}\cos2t+\frac12e^{-t}\sin2t.
\]
Therefore,
\[
\boxed{
x(t)=e^{-t}\left(\cos2t+\frac12\sin2t\right).
}
\]
The characteristic roots are $\lambda=-1\pm2i$,
so this agrees with the expected form:
an exponential decay multiplied by a wave.
\end{example}

% ------

\subsection{Completing the Square}

In Example 6, we rewrote the denominator to reveal
the shift. In general, we complete the square by
taking half the coefficient of $s$, squaring it,
and adding and subtracting that quantity:
\begin{align*}
s^2+bs+c
&=\left(s^2+bs+\frac{b^2}{4}\right)
+c-\frac{b^2}{4}\\
&=\boxed{\left(s+\frac{b}{2}\right)^2
+c-\frac{b^2}{4}.}
\end{align*}
For an irreducible quadratic, $c-b^2/4>0$.
This gives the form $(s-a)^2+\beta^2$ needed for
the shifted sine and cosine formulas, where
\[
a=-\frac{b}{2},
\qquad
\beta=\sqrt{c-\frac{b^2}{4}}.
\]
If the coefficient of $s^2$ is not $1$, factor it
out before completing the square.


\begin{example}[7]{
Find the inverse transform by completing the square
\[
\mathcal{L}^{-1}\left[\frac{7}{s^2+2s+50}\right].
\]
}
The denominator is an irreducible quadratic. Complete
the square to reveal the shift:
\[
s^2+2s+50=(s+1)^2+49.
\]
Thus,
\[
X(s)=\frac{7}{(s+1)^2+7^2}.
\]
The base transform is
\[
F(s)=\frac{7}{s^2+7^2},
\qquad
f(t)=\sin(7t).
\]
Since $X(s)=F(s+1)=F(s-(-1))$, the shift is $a=-1$.
Therefore,
\[
\boxed{
x(t)=\mathcal{L}^{-1}[X(s)]=e^{-t}\sin(7t).
}
\]
\end{example}
\section{Choosing the Right Formula}

We can now solve homogeneous second-order IVPs with
all four root types from our opening checklist.
The roots of the characteristic equation
\[
a\lambda^2+b\lambda+c=0
\]
help us anticipate the form of $X(s)$ and choose
the inverse transform.

\subsection{Distinct Real Roots}
For roots $\lambda_1\ne\lambda_2$, your transformed solution should be expressed in the form
\[
X(s)=\frac{A}{s-\lambda_1}+\frac{B}{s-\lambda_2}.
\]
This will require \textbf{partial fraction decomposition}. Each fraction produces an exponential:
\[
x(t)=Ae^{\lambda_1t}+Be^{\lambda_2t}.
\]
The shift formula is not required for this case.

\subsection{Repeated Real Root}
For a repeated root $\lambda$, your transformed solution should be expressed in the form
\[
X(s)=\frac{A}{s-\lambda}+\frac{B}{(s-\lambda)^2}.
\]
To get your solution in this form you may need \textbf{partial fraction decomposition}. Then you use the \textbf{shift formula} on the second term to get the solution 
\[
x(t)=Ae^{\lambda t}+Bte^{\lambda t}.
\]

\subsection{Complex Roots}
For roots $\alpha\pm i\beta$, you should aim to write your transformed solution in the form
\[
X(s)=
A\frac{s-\alpha}{(s-\alpha)^2+\beta^2}
+B\frac{\beta}{(s-\alpha)^2+\beta^2}.
\]
You will need to \textbf{complete the square}. Then apply the \textbf{shift formula} to obtain:
\[
x(t)=e^{\alpha t}\bigl(A\cos\beta t+B\sin\beta t\bigr).
\]
Purely imaginary roots are the special case $\alpha=0$,
so no exponential shift is needed.



\end{document}